% Generated draft: requires mathematical and source review.
\documentclass[11pt]{article}
\usepackage[margin=25mm]{geometry}
\usepackage{fontspec}
\setmainfont{Latin Modern Roman}
\usepackage{amsmath,amssymb,mathtools,unicode-math}
\setmathfont{Latin Modern Math}
\usepackage{enumitem,needspace,xurl,hyperref}
\setlength{\parindent}{0pt}
\setlength{\parskip}{5pt}
\setlength{\emergencystretch}{3em}
\newcommand{\N}{\mathbb{N}}
\newcommand{\Z}{\mathbb{Z}}
\newcommand{\Q}{\mathbb{Q}}
\newcommand{\R}{\mathbb{R}}
\newcommand{\C}{\mathbb{C}}
\newcommand{\F}{\mathbb{F}}
\newcommand{\A}{\mathbb{A}}
\newcommand{\PP}{\mathbb{P}}
\newcommand{\E}{\mathbb{E}}
\newcommand{\eps}{\varepsilon}
\newcommand{\dd}{\,\mathrm{d}}
\begin{document}
\section*{UnsolvedMath ID: 701578}
\textbf{Source problem code:} catalogue entry
\subsection{Definitions and mathematical statement}
For a finite group \(H\), let \(\operatorname{Irr}(H)\) denote its irreducible complex characters. An element \(h\in H\) is \emph{vanishing} if \(\chi(h)=0\) for some \(\chi\in\operatorname{Irr}(H)\). Define the set of orders of vanishing elements by
\[
\operatorname{Vo}(H)=\{\operatorname{ord}(h):h\in H,\ \exists\chi\in\operatorname{Irr}(H)\text{ with }\chi(h)=0\},
\]
where \(\operatorname{ord}(h)\) is the least positive integer \(n\) such that \(h^n=1_H\). This is a set, so multiplicities are disregarded. A simple group means a nontrivial group whose only normal subgroups are the identity subgroup and the whole group; the supplied statement imposes no nonabelian restriction.

Is it true that, for every finite group \(G\) and every finite simple group \(S\),
\[
\bigl(|G|=|S|\ \text{and}\ \operatorname{Vo}(G)=\operatorname{Vo}(S)\bigr)
\quad\Longrightarrow\quad G\cong S?
\]
\subsection{Short English statement}
Must a finite group be isomorphic to a finite simple group whenever they have the same cardinality and exactly the same set of orders of elements on which some irreducible complex character vanishes?
\subsection{Sources}
\begin{itemize}
\item Local UnsolvedMath registry, permanent ID \(701578\), source code \texttt{catalogue entry}, titled \emph{Kourovka Notebook Problem 19.30}. The supplied registry labels the problem open; this label is not independently verified here.
\item \emph{The Kourovka Notebook (Unsolved Problems in Group Theory)}, original problem 19.30, attributed in the supplied record to M. Foroudi Ghasemabadi and A. Iranmanesh, proposed in 2018. The supplied background identifies the 19th issue (2018) and the 21st issue edition (2026), with identifier \texttt{arXiv:1401.0300v48}.
\item The supplied collection description credits E. I. Khukhro and V. D. Mazurov as editors of the Kourovka Notebook collection covering Issues 1--21 (1965--2026).
\end{itemize}
\end{document}
